\section{Computational methods}\label{app:code}
\begin{itemize}
\item All computations were done with the author's Python code, which, together with the computed data, is available from the author. We implemented Koornwinder polynomials and the polynomial representations of the DAHA of types $D_N$, $C_N$ and $GL_N$; rational function arithmetic uses FLINT. Values at general $N$ were obtained by interpolating fixed-$N$ values as Laurent polynomials in $t^N$ and checking at the $N$ not used (\S\ref{sec:generalN}). For invariants with many terms, the linear system was solved degree by degree in $q^{1/2}$ with the range of powers restricted. For torus knots the heaviest computation is the three-box $T(2,9)$ on the $D_{14}$ DAHA, 90 minutes and 16\,GB per item.
\item For iterated torus knots (\S\ref{sec:cable}) the polynomial from the pattern step is expanded in the basis values $\{\hat\gamma_{\mathrm{comp}}(H_\nu)\}$ of the companion step, which are computed with coefficients kept as Laurent polynomials without cancellation. The heaviest item is one basis value of the $(2,13)$-cable on $D_{15}$, about 2 hours and 19\,GB.
\item Every Result in the text is checked by an automated test. The checks include the reproduction of the values of \cite{Marino09}, agreement with Kauffman polynomials at $t = q$, agreement with direct fixed-$N$ computations, agreement with the invariants of \cite{KN17}, and the symmetry $R \leftrightarrow R^T$.
\item The formulas in Appendices \ref{app:torus}--\ref{app:cover} and \ref{app:unref} were generated automatically from the computed data, not transcribed by hand.
\end{itemize}

\section{Torus knots in the fundamental representation}\label{app:torus}
$g^{\mathrm{true}}_{\yone}$ in bold variables, the $G$-expansion of the $c = 1$ part, and $\Pi_m$ of the $c = 2$ part.
$m = 1$ ($T(2,3)$, 11 terms): coefficients of $g^{\mathrm{true}}_{\yone}$ by powers of $\ba$
\begin{align*}
[\ba^{-6}]\colon\ & \bt^{-4} \\
[\ba^{-5}]\colon\ & \bq \bt^{-2} -\bq \bt^{-4} \\
[\ba^{-4}]\colon\ & -\bq^{2} \bt^{-2} -\bt^{-2} -\bq^{-2} \bt^{-4} \\
[\ba^{-3}]\colon\ & -\bq^{-1} \bt^{-2} +\bq^{-1} \bt^{-4} \\
[\ba^{-2}]\colon\ & \bq^{2} +\bt^{-2} +\bq^{-2} \bt^{-2}
\end{align*}

$m = 1$: $G$-expansion of the $c = 1$ part (per $\lambda^\beta$, $\qr = (q/t)^{1/2}$)
\begin{align*}
[\lambda^{-6}]\colon\ & \qr^{-4} \\
[\lambda^{-4}]\colon\ & -G \qr^{-3} -2 \qr^{-2} -\qr^{-4} \\
[\lambda^{-2}]\colon\ & G \qr^{-1} +1 +2 \qr^{-2}
\end{align*}

$m = 1$: the same, factorized per $\lambda^\beta$ ($\qr^2 - 1 = \bt^2 - 1$)
\begin{align*}
[\lambda^{-6}]\colon\ & \qr^{-4} \\
[\lambda^{-4}]\colon\ & \qr^{-4}\,\bigl(- G \qr - 2 \qr^{2} - 1\bigr) \\
[\lambda^{-2}]\colon\ & \qr^{-2}\,\bigl(G \qr + \qr^{2} + 2\bigr)
\end{align*}

\begin{align*}
\Pi_{1} &= 1
\end{align*}

$m = 2$ ($T(2,5)$, 27 terms): coefficients of $g^{\mathrm{true}}_{\yone}$ by powers of $\ba$
\begin{align*}
[\ba^{-10}]\colon\ & \bt^{-8} \\
[\ba^{-9}]\colon\ & \bq \bt^{-6} -\bq \bt^{-8} \\
[\ba^{-8}]\colon\ & \bq^{2} \bt^{-4} -\bq^{2} \bt^{-6} +\bq^{-2} \bt^{-6} -\bq^{-2} \bt^{-8} \\
[\ba^{-7}]\colon\ & \bq^{3} \bt^{-2} -\bq^{3} \bt^{-4} +\bq^{-1} \bt^{-4} -2 \bq^{-1} \bt^{-6} +\bq^{-1} \bt^{-8} \\
[\ba^{-6}]\colon\ & -\bq^{4} \bt^{-2} -\bq^{2} \bt^{-2} -2 \bt^{-4} +\bt^{-6} -\bq^{-2} \bt^{-4} -\bq^{-4} \bt^{-6} \\
[\ba^{-5}]\colon\ & -\bq \bt^{-2} +\bq \bt^{-4} -\bq^{-3} \bt^{-4} +\bq^{-3} \bt^{-6} \\
[\ba^{-4}]\colon\ & \bq^{4} +\bq^{2} \bt^{-2} +\bt^{-2} +\bq^{-2} \bt^{-4} +\bq^{-4} \bt^{-4}
\end{align*}

$m = 2$: $G$-expansion of the $c = 1$ part (per $\lambda^\beta$, $\qr = (q/t)^{1/2}$)
\begin{align*}
[\lambda^{-10}]\colon\ & \qr^{-8} \\
[\lambda^{-8}]\colon\ & G \qr^{-5} -G \qr^{-7} +\qr^{-4} -\qr^{-8} \\
[\lambda^{-6}]\colon\ & -G^{2} \qr^{-4} -3 G \qr^{-3} -2 G \qr^{-5} -2 \qr^{-2} -3 \qr^{-4} \\
[\lambda^{-4}]\colon\ & G^{2} \qr^{-2} +2 G \qr^{-1} +3 G \qr^{-3} +1 +2 \qr^{-2} +2 \qr^{-4}
\end{align*}

$m = 2$: the same, factorized per $\lambda^\beta$ ($\qr^2 - 1 = \bt^2 - 1$)
\begin{align*}
[\lambda^{-10}]\colon\ & \qr^{-8} \\
[\lambda^{-8}]\colon\ & \qr^{-8}\,\bigl((\qr^2 - 1)(G \qr + \qr^{2} + 1)\bigr) \\
[\lambda^{-6}]\colon\ & \text{irreducible (as above)} \\
[\lambda^{-4}]\colon\ & \text{irreducible (as above)}
\end{align*}

\begin{align*}
\Pi_{2} &= \ba^{2} \bq^{4} \bt^{4} +\ba^{2} \bt^{2} +\bq^{2}
\end{align*}

$m = 3$ ($T(2,7)$, 49 terms): coefficients of $g^{\mathrm{true}}_{\yone}$ by powers of $\ba$
\begin{align*}
[\ba^{-14}]\colon\ & \bt^{-12} \\
[\ba^{-13}]\colon\ & \bq \bt^{-10} -\bq \bt^{-12} \\
[\ba^{-12}]\colon\ & \bq^{2} \bt^{-8} -\bq^{2} \bt^{-10} +\bq^{-2} \bt^{-10} -\bq^{-2} \bt^{-12} \\
[\ba^{-11}]\colon\ & \bq^{3} \bt^{-6} -\bq^{3} \bt^{-8} +\bq^{-1} \bt^{-8} -2 \bq^{-1} \bt^{-10} +\bq^{-1} \bt^{-12} \\
[\ba^{-10}]\colon\ & \bq^{4} \bt^{-4} -\bq^{4} \bt^{-6} +\bt^{-6} -2 \bt^{-8} +\bt^{-10} +\bq^{-4} \bt^{-8} \\
& \quad -\bq^{-4} \bt^{-10} \\
[\ba^{-9}]\colon\ & \bq^{5} \bt^{-2} -\bq^{5} \bt^{-4} +\bq \bt^{-4} -2 \bq \bt^{-6} +\bq \bt^{-8} +\bq^{-3} \bt^{-6} \\
& \quad -2 \bq^{-3} \bt^{-8} +\bq^{-3} \bt^{-10} \\
[\ba^{-8}]\colon\ & -\bq^{6} \bt^{-2} -\bq^{4} \bt^{-2} -2 \bq^{2} \bt^{-4} +\bq^{2} \bt^{-6} -\bt^{-4} -2 \bq^{-2} \bt^{-6} \\
& \quad +\bq^{-2} \bt^{-8} -\bq^{-4} \bt^{-6} -\bq^{-6} \bt^{-8} \\
[\ba^{-7}]\colon\ & -\bq^{3} \bt^{-2} +\bq^{3} \bt^{-4} -\bq^{-1} \bt^{-4} +\bq^{-1} \bt^{-6} -\bq^{-5} \bt^{-6} +\bq^{-5} \bt^{-8} \\
[\ba^{-6}]\colon\ & \bq^{6} +\bq^{4} \bt^{-2} +\bq^{2} \bt^{-2} +\bt^{-4} +\bq^{-2} \bt^{-4} +\bq^{-4} \bt^{-6} \\
& \quad +\bq^{-6} \bt^{-6}
\end{align*}

$m = 3$: $G$-expansion of the $c = 1$ part (per $\lambda^\beta$, $\qr = (q/t)^{1/2}$)
\begin{align*}
[\lambda^{-14}]\colon\ & \qr^{-12} \\
[\lambda^{-12}]\colon\ & G \qr^{-9} -G \qr^{-11} +\qr^{-8} -\qr^{-12} \\
[\lambda^{-10}]\colon\ & G^{2} \qr^{-6} -G^{2} \qr^{-8} +2 G \qr^{-5} -2 G \qr^{-9} +\qr^{-4} -\qr^{-8} \\
[\lambda^{-8}]\colon\ & -G^{3} \qr^{-5} -4 G^{2} \qr^{-4} -3 G^{2} \qr^{-6} -5 G \qr^{-3} -7 G \qr^{-5} -2 G \qr^{-7} \\
& \quad -2 \qr^{-2} -3 \qr^{-4} -2 \qr^{-6} \\
[\lambda^{-6}]\colon\ & G^{3} \qr^{-3} +3 G^{2} \qr^{-2} +4 G^{2} \qr^{-4} +3 G \qr^{-1} +6 G \qr^{-3} +5 G \qr^{-5} \\
& \quad +1 +2 \qr^{-2} +2 \qr^{-4} +2 \qr^{-6}
\end{align*}

$m = 3$: the same, factorized per $\lambda^\beta$ ($\qr^2 - 1 = \bt^2 - 1$)
\begin{align*}
[\lambda^{-14}]\colon\ & \qr^{-12} \\
[\lambda^{-12}]\colon\ & \qr^{-12}\,\bigl((\qr^2 - 1)(G \qr + \qr^{2} + 1)\bigr) \\
[\lambda^{-10}]\colon\ & \qr^{-9}\,\bigl((\qr^2 - 1)(G^{2} \qr + 2 G \qr^{2} + 2 G + \qr^{3} + \qr)\bigr) \\
[\lambda^{-8}]\colon\ & \text{irreducible (as above)} \\
[\lambda^{-6}]\colon\ & \text{irreducible (as above)}
\end{align*}

\begin{align*}
\Pi_{3} &= \ba^{4} \bq^{8} \bt^{8} +\ba^{4} \bq^{4} \bt^{6} +\ba^{4} \bt^{4} +\ba^{2} \bq^{6} \bt^{4} +\ba^{2} \bq^{2} \bt^{2} +\bq^{4}
\end{align*}


\section{DAHA: reduced invariants at general \texorpdfstring{$N$}{N}}\label{app:daha}
$\Rred_R(t^N) = T_R^{-6}\times$ mirror$(Z)/K_R(t^\rho)$ for the trefoil in the uniform normalization, where $Z$ is the unnormalized DAHA-Jones value. The powers of $q, t$ can be half-integral.
$\Rred_{\yone}(t^N)$ (coefficients of $(t^N)^j$):
\begin{align*}
[(t^N)^{-5}]\colon\ & -q^{-2} t^{5} +q^{-2} t^{4} \\
[(t^N)^{-4}]\colon\ & -q^{-1} t^{4} +q^{-2} t^{4} -q^{-2} t^{3} \\
[(t^N)^{-3}]\colon\ & q^{-1} t^{3} -q^{-1} t^{2} \\
[(t^N)^{-2}]\colon\ & t^{2} +q^{-1} t
\end{align*}

$\Rred_{\ytwo}(t^N)$ (coefficients of $(t^N)^j$):
\begin{align*}
[(t^N)^{-10}]\colon\ & q^{-9} t^{10} -q^{-9} t^{9} -q^{-10} t^{9} +q^{-10} t^{8} \\
[(t^N)^{-9}]\colon\ & q^{-7} t^{9} -q^{-7} t^{8} -q^{-9} t^{9} +2 q^{-9} t^{8} -q^{-9} t^{7} +q^{-10} t^{8} \\
& \quad -q^{-10} t^{7} \\
[(t^N)^{-8}]\colon\ & q^{-5} t^{8} -q^{-6} t^{8} +q^{-6} t^{7} -2 q^{-7} t^{8} +3 q^{-7} t^{7} -q^{-7} t^{6} \\
& \quad +q^{-8} t^{7} -q^{-8} t^{6} -q^{-9} t^{7} +q^{-9} t^{6} \\
[(t^N)^{-7}]\colon\ & -q^{-4} t^{7} +q^{-4} t^{6} -2 q^{-5} t^{7} +2 q^{-5} t^{6} -q^{-6} t^{6} +q^{-6} t^{5} \\
& \quad +q^{-7} t^{7} -3 q^{-7} t^{6} +2 q^{-7} t^{5} -q^{-8} t^{6} +q^{-8} t^{5} \\
[(t^N)^{-6}]\colon\ & -q^{-2} t^{6} -q^{-3} t^{6} +q^{-4} t^{6} -2 q^{-4} t^{5} +q^{-5} t^{6} -3 q^{-5} t^{5} \\
& \quad +q^{-5} t^{4} +q^{-6} t^{6} -q^{-6} t^{4} +q^{-7} t^{5} -q^{-7} t^{4} \\
[(t^N)^{-5}]\colon\ & q^{-2} t^{5} -q^{-2} t^{4} +q^{-3} t^{5} -q^{-3} t^{4} +q^{-4} t^{4} -q^{-4} t^{3} \\
& \quad +q^{-5} t^{4} -q^{-5} t^{3} \\
[(t^N)^{-4}]\colon\ & t^{4} +q^{-2} t^{3} +q^{-3} t^{3} +q^{-4} t^{2}
\end{align*}

$\Rred_{\yoneone}(t^N)$ (coefficients of $(t^N)^j$):
\begin{align*}
[(t^N)^{-10}]\colon\ & q^{-4} t^{16} -q^{-4} t^{15} -q^{-5} t^{15} +q^{-5} t^{14} \\
[(t^N)^{-9}]\colon\ & q^{-3} t^{15} -q^{-3} t^{13} -q^{-4} t^{14} +q^{-4} t^{13} +q^{-4} t^{12} -q^{-4} t^{11} \\
[(t^N)^{-8}]\colon\ & q^{-2} t^{13} -2 q^{-3} t^{13} +2 q^{-3} t^{11} +q^{-4} t^{12} -q^{-4} t^{11} -q^{-4} t^{10} \\
& \quad +q^{-4} t^{9} \\
[(t^N)^{-7}]\colon\ & -q^{-2} t^{12} -q^{-2} t^{11} +q^{-2} t^{10} +q^{-2} t^{9} +q^{-3} t^{11} -q^{-3} t^{10} \\
& \quad -2 q^{-3} t^{9} +q^{-3} t^{8} +q^{-3} t^{7} \\
[(t^N)^{-6}]\colon\ & -q^{-1} t^{10} -q^{-1} t^{9} +q^{-2} t^{10} +q^{-2} t^{9} -3 q^{-2} t^{8} -q^{-2} t^{7} \\
& \quad +q^{-3} t^{9} +q^{-3} t^{7} -q^{-3} t^{6} -q^{-3} t^{5} \\
[(t^N)^{-5}]\colon\ & q^{-1} t^{8} -q^{-1} t^{6} +q^{-2} t^{6} -q^{-2} t^{4} \\
[(t^N)^{-4}]\colon\ & t^{6} +q^{-1} t^{5} +q^{-1} t^{4} +q^{-2} t^{2}
\end{align*}

$\Rred_{\mathrm{adj}}$ of the adjoint $(1;1)$ of the cover (coefficients of powers of $\lambda = t^{N_A/2}$):
\begin{align*}
[\lambda^{-5}]\colon\ & q^{-4} t^{5} -q^{-4} t^{4} -q^{-5} t^{5} +q^{-5} t^{4} \\
[\lambda^{-9/2}]\colon\ & 0 \\
[\lambda^{-4}]\colon\ & q^{-2} t^{4} -2 q^{-3} t^{4} +2 q^{-3} t^{3} +q^{-4} t^{4} -2 q^{-4} t^{3} +q^{-4} t^{2} \\
[\lambda^{-7/2}]\colon\ & 0 \\
[\lambda^{-3}]\colon\ & -2 q^{-1} t^{3} +q^{-2} t^{3} -3 q^{-2} t^{2} +q^{-3} t^{3} +q^{-3} t^{2} -2 q^{-3} t \\
[\lambda^{-5/2}]\colon\ & 0 \\
[\lambda^{-2}]\colon\ & t^{2} +2 q^{-1} t +q^{-2}
\end{align*}


\section{Two-box crosscap part at large \texorpdfstring{$N$}{N}}\label{app:largeN}
$\hat g_{\ytwo} + \hat g_{\yoneone}$, $(1+t)(\hat g_{\ytwo} - \hat g_{\yoneone})$ and the quotient by the universal factor (bold variables) for the trefoil. These document the integrality (Result \ref{cr:largeN}) and the universal factor (\S\ref{sec:universal}); the signs per power of $\ba$ do not align (Result \ref{cr:signs}). The comparison of sectors with and without visible positivity is in Appendix \ref{app:cover}.
$\hat g_{\ytwo} + \hat g_{\yoneone}$ (integer Laurent, 262 terms):
\begin{align*}
[\ba^{-12}]\colon\ & -\bq^{11} \bt^{-9} -\bq^{9} \bt^{-9} -\bq^{7} \bt^{-9} -\bq^{5} \bt^{-9} -\bq^{3} \bt^{-9} -\bq \bt^{-9} \\
& \quad +\bq^{-1} \bt^{-9} +\bq^{-3} \bt^{-11} +\bq^{-5} \bt^{-13} +\bq^{-7} \bt^{-15} +\bq^{-9} \bt^{-17} +\bq^{-11} \bt^{-19} \\
[\ba^{-11}]\colon\ & -\bq^{12} \bt^{-7} +\bq^{12} \bt^{-9} -2 \bq^{10} \bt^{-7} +\bq^{10} \bt^{-9} -2 \bq^{8} \bt^{-7} -2 \bq^{6} \bt^{-7} \\
& \quad -\bq^{6} \bt^{-9} -2 \bq^{4} \bt^{-7} -2 \bq^{4} \bt^{-9} -2 \bq^{2} \bt^{-7} -2 \bq^{2} \bt^{-9} -\bq^{2} \bt^{-11} \\
& \quad -4 \bt^{-9} -2 \bt^{-11} -4 \bq^{-2} \bt^{-11} -\bq^{-2} \bt^{-13} -4 \bq^{-4} \bt^{-13} -3 \bq^{-6} \bt^{-15} \\
& \quad -2 \bq^{-8} \bt^{-17} -\bq^{-10} \bt^{-19} \\
[\ba^{-10}]\colon\ & \bq^{13} \bt^{-7} -\bq^{11} \bt^{-5} +3 \bq^{11} \bt^{-7} -2 \bq^{9} \bt^{-5} +3 \bq^{9} \bt^{-7} +2 \bq^{9} \bt^{-9} \\
& \quad -3 \bq^{7} \bt^{-5} +3 \bq^{7} \bt^{-7} +3 \bq^{7} \bt^{-9} -4 \bq^{5} \bt^{-5} +2 \bq^{5} \bt^{-7} +5 \bq^{5} \bt^{-9} \\
& \quad -4 \bq^{3} \bt^{-5} +\bq^{3} \bt^{-7} +5 \bq^{3} \bt^{-9} +\bq^{3} \bt^{-11} -2 \bq \bt^{-5} -5 \bq \bt^{-7} \\
& \quad +5 \bq \bt^{-9} +3 \bq \bt^{-11} -4 \bq^{-1} \bt^{-7} -2 \bq^{-1} \bt^{-9} +4 \bq^{-1} \bt^{-11} +\bq^{-1} \bt^{-13} \\
& \quad -5 \bq^{-3} \bt^{-9} -2 \bq^{-3} \bt^{-11} +4 \bq^{-3} \bt^{-13} -4 \bq^{-5} \bt^{-11} -\bq^{-5} \bt^{-13} +2 \bq^{-5} \bt^{-15} \\
& \quad -4 \bq^{-7} \bt^{-13} +\bq^{-7} \bt^{-17} -3 \bq^{-9} \bt^{-15} -2 \bq^{-11} \bt^{-17} -\bq^{-13} \bt^{-19} \\
[\ba^{-9}]\colon\ & 2 \bq^{12} \bt^{-5} -\bq^{12} \bt^{-7} +4 \bq^{10} \bt^{-5} -\bq^{10} \bt^{-9} +5 \bq^{8} \bt^{-5} +2 \bq^{8} \bt^{-7} \\
& \quad -\bq^{8} \bt^{-9} +6 \bq^{6} \bt^{-5} +5 \bq^{6} \bt^{-7} -\bq^{6} \bt^{-9} +6 \bq^{4} \bt^{-5} +9 \bq^{4} \bt^{-7} \\
& \quad +4 \bq^{2} \bt^{-5} +13 \bq^{2} \bt^{-7} +4 \bq^{2} \bt^{-9} -\bq^{2} \bt^{-11} +12 \bt^{-7} +10 \bt^{-9} \\
& \quad +15 \bq^{-2} \bt^{-9} +6 \bq^{-2} \bt^{-11} -\bq^{-2} \bt^{-13} +13 \bq^{-4} \bt^{-11} +2 \bq^{-4} \bt^{-13} +10 \bq^{-6} \bt^{-13} \\
& \quad +6 \bq^{-8} \bt^{-15} +3 \bq^{-10} \bt^{-17} +\bq^{-12} \bt^{-19} \\
[\ba^{-8}]\colon\ & -\bq^{13} \bt^{-5} +\bq^{11} \bt^{-3} -2 \bq^{11} \bt^{-5} -\bq^{11} \bt^{-7} +2 \bq^{9} \bt^{-3} -\bq^{9} \bt^{-5} \\
& \quad -4 \bq^{9} \bt^{-7} +4 \bq^{7} \bt^{-3} -6 \bq^{7} \bt^{-7} -\bq^{7} \bt^{-9} +5 \bq^{5} \bt^{-3} +4 \bq^{5} \bt^{-5} \\
& \quad -10 \bq^{5} \bt^{-7} -2 \bq^{5} \bt^{-9} +4 \bq^{3} \bt^{-3} +10 \bq^{3} \bt^{-5} -11 \bq^{3} \bt^{-7} -6 \bq^{3} \bt^{-9} \\
& \quad +17 \bq \bt^{-5} -7 \bq \bt^{-7} -10 \bq \bt^{-9} -\bq \bt^{-11} +5 \bq^{-1} \bt^{-5} +14 \bq^{-1} \bt^{-7} \\
& \quad -13 \bq^{-1} \bt^{-9} -5 \bq^{-1} \bt^{-11} +8 \bq^{-3} \bt^{-7} +7 \bq^{-3} \bt^{-9} -11 \bq^{-3} \bt^{-11} -\bq^{-3} \bt^{-13} \\
& \quad +7 \bq^{-5} \bt^{-9} +3 \bq^{-5} \bt^{-11} -7 \bq^{-5} \bt^{-13} +5 \bq^{-7} \bt^{-11} +\bq^{-7} \bt^{-13} -3 \bq^{-7} \bt^{-15} \\
& \quad +4 \bq^{-9} \bt^{-13} -\bq^{-9} \bt^{-17} +2 \bq^{-11} \bt^{-15} +\bq^{-13} \bt^{-17} \\
[\ba^{-7}]\colon\ & -\bq^{12} \bt^{-3} -2 \bq^{10} \bt^{-3} -2 \bq^{10} \bt^{-5} +\bq^{10} \bt^{-7} -4 \bq^{8} \bt^{-3} -4 \bq^{8} \bt^{-5} \\
& \quad +\bq^{8} \bt^{-7} -6 \bq^{6} \bt^{-3} -8 \bq^{6} \bt^{-5} +\bq^{6} \bt^{-7} -6 \bq^{4} \bt^{-3} -15 \bq^{4} \bt^{-5} \\
& \quad -\bq^{4} \bt^{-7} +\bq^{4} \bt^{-9} -2 \bq^{2} \bt^{-3} -22 \bq^{2} \bt^{-5} -7 \bq^{2} \bt^{-7} +2 \bq^{2} \bt^{-9} \\
& \quad -11 \bt^{-5} -22 \bt^{-7} +\bt^{-11} -19 \bq^{-2} \bt^{-7} -12 \bq^{-2} \bt^{-9} +2 \bq^{-2} \bt^{-11} \\
& \quad -17 \bq^{-4} \bt^{-9} -5 \bq^{-4} \bt^{-11} +\bq^{-4} \bt^{-13} -12 \bq^{-6} \bt^{-11} -\bq^{-6} \bt^{-13} -7 \bq^{-8} \bt^{-13} \\
& \quad -3 \bq^{-10} \bt^{-15} -\bq^{-12} \bt^{-17} \\
[\ba^{-6}]\colon\ & \bq^{11} \bt^{-5} -\bq^{9} \bt^{-3} +2 \bq^{9} \bt^{-5} -\bq^{7} \bt^{-1} -2 \bq^{7} \bt^{-3} +3 \bq^{7} \bt^{-5} \\
& \quad +\bq^{7} \bt^{-7} -\bq^{5} \bt^{-1} -6 \bq^{5} \bt^{-3} +5 \bq^{5} \bt^{-5} +3 \bq^{5} \bt^{-7} -11 \bq^{3} \bt^{-3} \\
& \quad +4 \bq^{3} \bt^{-5} +8 \bq^{3} \bt^{-7} -7 \bq \bt^{-3} -7 \bq \bt^{-5} +13 \bq \bt^{-7} +2 \bq \bt^{-9} \\
& \quad -19 \bq^{-1} \bt^{-5} +11 \bq^{-1} \bt^{-7} +7 \bq^{-1} \bt^{-9} -4 \bq^{-3} \bt^{-5} -11 \bq^{-3} \bt^{-7} +12 \bq^{-3} \bt^{-9} \\
& \quad +2 \bq^{-3} \bt^{-11} -5 \bq^{-5} \bt^{-7} -4 \bq^{-5} \bt^{-9} +8 \bq^{-5} \bt^{-11} -3 \bq^{-7} \bt^{-9} -\bq^{-7} \bt^{-11} \\
& \quad +3 \bq^{-7} \bt^{-13} -2 \bq^{-9} \bt^{-11} +\bq^{-9} \bt^{-15} -\bq^{-11} \bt^{-13} \\
[\ba^{-5}]\colon\ & \bq^{10} \bt^{-3} +\bq^{8} \bt^{-1} +2 \bq^{8} \bt^{-3} +2 \bq^{6} \bt^{-1} +5 \bq^{6} \bt^{-3} +2 \bq^{4} \bt^{-1} \\
& \quad +10 \bq^{4} \bt^{-3} +2 \bq^{4} \bt^{-5} -\bq^{4} \bt^{-7} +12 \bq^{2} \bt^{-3} +9 \bq^{2} \bt^{-5} -2 \bq^{2} \bt^{-7} \\
& \quad +3 \bt^{-3} +21 \bt^{-5} -\bt^{-7} -\bt^{-9} +9 \bq^{-2} \bt^{-5} +12 \bq^{-2} \bt^{-7} \\
& \quad -2 \bq^{-2} \bt^{-9} +10 \bq^{-4} \bt^{-7} +4 \bq^{-4} \bt^{-9} -\bq^{-4} \bt^{-11} +6 \bq^{-6} \bt^{-9} +\bq^{-6} \bt^{-11} \\
& \quad +3 \bq^{-8} \bt^{-11} +\bq^{-10} \bt^{-13} \\
[\ba^{-4}]\colon\ & \bq^{5} \bt^{-1} -\bq^{5} \bt^{-5} +\bq^{3} \bt^{-1} +2 \bq^{3} \bt^{-3} -3 \bq^{3} \bt^{-5} +7 \bq \bt^{-3} \\
& \quad -6 \bq \bt^{-5} -\bq \bt^{-7} +3 \bq^{-1} \bt^{-3} +\bq^{-1} \bt^{-5} -4 \bq^{-1} \bt^{-7} +6 \bq^{-3} \bt^{-5} \\
& \quad -5 \bq^{-3} \bt^{-7} -\bq^{-3} \bt^{-9} +\bq^{-5} \bt^{-5} +2 \bq^{-5} \bt^{-7} -3 \bq^{-5} \bt^{-9} +\bq^{-7} \bt^{-7} \\
& \quad -\bq^{-7} \bt^{-11} \\
[\ba^{-3}]\colon\ & -\bq^{6} \bt^{-1} -2 \bq^{4} \bt^{-1} -\bq^{4} \bt^{-3} -\bq^{2} \bt^{-1} -5 \bq^{2} \bt^{-3} +\bq^{2} \bt^{-5} \\
& \quad -6 \bt^{-3} -\bq^{-2} \bt^{-3} -5 \bq^{-2} \bt^{-5} +\bq^{-2} \bt^{-7} -2 \bq^{-4} \bt^{-5} -\bq^{-4} \bt^{-7} \\
& \quad -\bq^{-6} \bt^{-7} \\
[\ba^{-2}]\colon\ & -\bq^{-1} \bt^{-3} +\bq^{-1} \bt^{-5} \\
[\ba^{-1}]\colon\ & -\bt^{-1} +\bt^{-3}
\end{align*}

$(1 + t)(\hat g_{\ytwo} - \hat g_{\yoneone})$ (integer Laurent, 290 terms):
\begin{align*}
[\ba^{-12}]\colon\ & -\bq^{13} \bt^{-9} -\bq \bt^{-9} -\bq^{-1} \bt^{-9} +\bq^{-1} \bt^{-11} +\bq^{-3} \bt^{-11} -\bq^{-3} \bt^{-13} \\
& \quad -\bq^{-5} \bt^{-13} +\bq^{-5} \bt^{-15} +\bq^{-7} \bt^{-15} -\bq^{-7} \bt^{-17} -\bq^{-9} \bt^{-17} +\bq^{-9} \bt^{-19} \\
& \quad +\bq^{-11} \bt^{-19} \\
[\ba^{-11}]\colon\ & -\bq^{14} \bt^{-7} +\bq^{14} \bt^{-9} -\bq^{12} \bt^{-7} -\bq^{10} \bt^{-9} -\bq^{8} \bt^{-9} -\bq^{6} \bt^{-9} \\
& \quad -\bq^{4} \bt^{-11} -\bq^{2} \bt^{-11} -\bt^{-13} -\bq^{-2} \bt^{-13} -\bq^{-4} \bt^{-15} -\bq^{-6} \bt^{-15} \\
& \quad -\bq^{-8} \bt^{-19} -\bq^{-10} \bt^{-19} \\
[\ba^{-10}]\colon\ & \bq^{15} \bt^{-7} -\bq^{13} \bt^{-5} +2 \bq^{13} \bt^{-7} -\bq^{11} \bt^{-5} +2 \bq^{11} \bt^{-9} -\bq^{9} \bt^{-5} \\
& \quad +\bq^{9} \bt^{-9} -\bq^{7} \bt^{-5} -\bq^{7} \bt^{-7} +2 \bq^{7} \bt^{-9} -\bq^{5} \bt^{-7} +\bq^{5} \bt^{-11} \\
& \quad -\bq^{3} \bt^{-5} -\bq^{3} \bt^{-7} +2 \bq^{3} \bt^{-11} -\bq \bt^{-5} +\bq \bt^{-7} -\bq \bt^{-11} \\
& \quad +\bq \bt^{-13} +\bq^{-1} \bt^{-7} -2 \bq^{-1} \bt^{-9} -2 \bq^{-1} \bt^{-11} +3 \bq^{-1} \bt^{-13} -\bq^{-3} \bt^{-9} \\
& \quad +\bq^{-3} \bt^{-13} -\bq^{-5} \bt^{-13} +\bq^{-5} \bt^{-17} -\bq^{-7} \bt^{-15} +\bq^{-7} \bt^{-17} -\bq^{-9} \bt^{-15} \\
& \quad -\bq^{-11} \bt^{-19} -\bq^{-13} \bt^{-19} \\
[\ba^{-9}]\colon\ & 2 \bq^{14} \bt^{-5} -\bq^{14} \bt^{-7} +2 \bq^{12} \bt^{-5} +\bq^{12} \bt^{-7} -\bq^{12} \bt^{-9} +\bq^{10} \bt^{-5} \\
& \quad +2 \bq^{10} \bt^{-7} +\bq^{8} \bt^{-5} +3 \bq^{8} \bt^{-7} +4 \bq^{6} \bt^{-7} +\bq^{6} \bt^{-9} +2 \bq^{4} \bt^{-7} \\
& \quad +4 \bq^{4} \bt^{-9} -\bq^{4} \bt^{-11} +\bq^{2} \bt^{-7} +4 \bq^{2} \bt^{-9} +\bq^{2} \bt^{-11} +3 \bt^{-9} \\
& \quad +4 \bt^{-11} -\bt^{-13} +\bq^{-2} \bt^{-9} +3 \bq^{-2} \bt^{-11} +\bq^{-2} \bt^{-13} +\bq^{-4} \bt^{-11} \\
& \quad +4 \bq^{-4} \bt^{-13} +2 \bq^{-6} \bt^{-13} +2 \bq^{-6} \bt^{-15} +2 \bq^{-8} \bt^{-15} +\bq^{-8} \bt^{-17} +\bq^{-10} \bt^{-17} \\
& \quad +\bq^{-10} \bt^{-19} +\bq^{-12} \bt^{-19} \\
[\ba^{-8}]\colon\ & -\bq^{15} \bt^{-5} +\bq^{13} \bt^{-3} -\bq^{13} \bt^{-5} -\bq^{13} \bt^{-7} +\bq^{11} \bt^{-3} +\bq^{11} \bt^{-5} \\
& \quad -3 \bq^{11} \bt^{-7} +2 \bq^{9} \bt^{-3} +\bq^{9} \bt^{-5} -2 \bq^{9} \bt^{-7} -\bq^{9} \bt^{-9} +\bq^{7} \bt^{-3} \\
& \quad +4 \bq^{7} \bt^{-5} -4 \bq^{7} \bt^{-7} -\bq^{7} \bt^{-9} +6 \bq^{5} \bt^{-5} -\bq^{5} \bt^{-7} -4 \bq^{5} \bt^{-9} \\
& \quad +4 \bq^{3} \bt^{-5} +\bq^{3} \bt^{-7} -4 \bq^{3} \bt^{-9} -\bq^{3} \bt^{-11} +2 \bq \bt^{-5} +4 \bq \bt^{-7} \\
& \quad -\bq \bt^{-9} -4 \bq \bt^{-11} +\bq^{-1} \bt^{-5} +\bq^{-1} \bt^{-7} +3 \bq^{-1} \bt^{-9} -4 \bq^{-1} \bt^{-11} \\
& \quad -\bq^{-1} \bt^{-13} -\bq^{-3} \bt^{-7} +6 \bq^{-3} \bt^{-9} -4 \bq^{-3} \bt^{-13} +3 \bq^{-5} \bt^{-9} -2 \bq^{-5} \bt^{-13} \\
& \quad -\bq^{-5} \bt^{-15} -\bq^{-7} \bt^{-11} +3 \bq^{-7} \bt^{-13} -\bq^{-7} \bt^{-15} -\bq^{-7} \bt^{-17} +2 \bq^{-9} \bt^{-13} \\
& \quad -\bq^{-9} \bt^{-17} +\bq^{-11} \bt^{-17} +\bq^{-13} \bt^{-17} \\
[\ba^{-7}]\colon\ & -\bq^{14} \bt^{-3} -\bq^{12} \bt^{-3} -2 \bq^{12} \bt^{-5} +\bq^{12} \bt^{-7} -2 \bq^{10} \bt^{-3} -2 \bq^{10} \bt^{-5} \\
& \quad -2 \bq^{8} \bt^{-3} -4 \bq^{8} \bt^{-5} -7 \bq^{6} \bt^{-5} -2 \bq^{6} \bt^{-7} +\bq^{6} \bt^{-9} -5 \bq^{4} \bt^{-5} \\
& \quad -6 \bq^{4} \bt^{-7} +\bq^{4} \bt^{-9} -\bq^{2} \bt^{-5} -9 \bq^{2} \bt^{-7} -2 \bq^{2} \bt^{-9} +\bq^{2} \bt^{-11} \\
& \quad +\bt^{-5} -7 \bt^{-7} -6 \bt^{-9} +\bt^{-11} -\bq^{-2} \bt^{-7} -7 \bq^{-2} \bt^{-9} \\
& \quad -3 \bq^{-2} \bt^{-11} +\bq^{-2} \bt^{-13} -3 \bq^{-4} \bt^{-9} -5 \bq^{-4} \bt^{-11} -2 \bq^{-6} \bt^{-11} -4 \bq^{-6} \bt^{-13} \\
& \quad -3 \bq^{-8} \bt^{-13} -\bq^{-8} \bt^{-15} -\bq^{-10} \bt^{-15} -\bq^{-10} \bt^{-17} -\bq^{-12} \bt^{-17} \\
[\ba^{-6}]\colon\ & \bq^{13} \bt^{-5} -\bq^{11} \bt^{-3} +\bq^{11} \bt^{-5} -\bq^{9} \bt^{-1} -\bq^{9} \bt^{-3} +\bq^{9} \bt^{-5} \\
& \quad +\bq^{9} \bt^{-7} -4 \bq^{7} \bt^{-3} +2 \bq^{7} \bt^{-5} +2 \bq^{7} \bt^{-7} -4 \bq^{5} \bt^{-3} -\bq^{5} \bt^{-5} \\
& \quad +5 \bq^{5} \bt^{-7} -7 \bq^{3} \bt^{-5} +5 \bq^{3} \bt^{-7} +2 \bq^{3} \bt^{-9} +\bq \bt^{-3} -9 \bq \bt^{-5} \\
& \quad +3 \bq \bt^{-7} +5 \bq \bt^{-9} -\bq^{-1} \bt^{-5} -4 \bq^{-1} \bt^{-7} +3 \bq^{-1} \bt^{-9} +2 \bq^{-1} \bt^{-11} \\
& \quad +\bq^{-3} \bt^{-5} -5 \bq^{-3} \bt^{-7} -2 \bq^{-3} \bt^{-9} +6 \bq^{-3} \bt^{-11} -\bq^{-5} \bt^{-7} -3 \bq^{-5} \bt^{-9} \\
& \quad +3 \bq^{-5} \bt^{-11} +\bq^{-5} \bt^{-13} -\bq^{-7} \bt^{-9} -\bq^{-7} \bt^{-11} +\bq^{-7} \bt^{-13} +\bq^{-7} \bt^{-15} \\
& \quad -\bq^{-9} \bt^{-13} +\bq^{-9} \bt^{-15} -\bq^{-11} \bt^{-13} \\
[\ba^{-5}]\colon\ & \bq^{12} \bt^{-3} +\bq^{10} \bt^{-1} +\bq^{10} \bt^{-3} +\bq^{8} \bt^{-1} +3 \bq^{8} \bt^{-3} +5 \bq^{6} \bt^{-3} \\
& \quad +2 \bq^{6} \bt^{-5} -\bq^{6} \bt^{-7} +2 \bq^{4} \bt^{-3} +7 \bq^{4} \bt^{-5} -\bq^{4} \bt^{-7} -\bq^{2} \bt^{-3} \\
& \quad +10 \bq^{2} \bt^{-5} +\bq^{2} \bt^{-7} -\bq^{2} \bt^{-9} -\bt^{-3} +4 \bt^{-5} +7 \bt^{-7} \\
& \quad -\bt^{-9} -\bq^{-2} \bt^{-5} +8 \bq^{-2} \bt^{-7} +2 \bq^{-2} \bt^{-9} -\bq^{-2} \bt^{-11} +2 \bq^{-4} \bt^{-7} \\
& \quad +4 \bq^{-4} \bt^{-9} +2 \bq^{-6} \bt^{-9} +2 \bq^{-6} \bt^{-11} +\bq^{-8} \bt^{-11} +\bq^{-8} \bt^{-13} +\bq^{-10} \bt^{-13} \\
[\ba^{-4}]\colon\ & \bq^{7} \bt^{-1} -\bq^{7} \bt^{-5} -\bq^{5} \bt^{-1} +2 \bq^{5} \bt^{-3} -2 \bq^{5} \bt^{-5} -\bq^{3} \bt^{-1} \\
& \quad +5 \bq^{3} \bt^{-3} -3 \bq^{3} \bt^{-5} -\bq^{3} \bt^{-7} +\bq \bt^{-3} +\bq \bt^{-5} -3 \bq \bt^{-7} \\
& \quad -\bq^{-1} \bt^{-3} +5 \bq^{-1} \bt^{-5} -3 \bq^{-1} \bt^{-7} -\bq^{-1} \bt^{-9} +\bq^{-3} \bt^{-7} -2 \bq^{-3} \bt^{-9} \\
& \quad -\bq^{-5} \bt^{-5} +3 \bq^{-5} \bt^{-7} -\bq^{-5} \bt^{-9} -\bq^{-5} \bt^{-11} +\bq^{-7} \bt^{-7} -\bq^{-7} \bt^{-11} \\
[\ba^{-3}]\colon\ & -\bq^{8} \bt^{-1} -\bq^{6} \bt^{-1} -\bq^{6} \bt^{-3} +\bq^{4} \bt^{-1} -4 \bq^{4} \bt^{-3} +\bq^{4} \bt^{-5} \\
& \quad +\bq^{2} \bt^{-1} -3 \bq^{2} \bt^{-3} -\bq^{2} \bt^{-5} +\bt^{-3} -5 \bt^{-5} +\bt^{-7} \\
& \quad +\bq^{-2} \bt^{-3} -3 \bq^{-2} \bt^{-5} -2 \bq^{-4} \bt^{-7} -\bq^{-6} \bt^{-7} \\
[\ba^{-2}]\colon\ & \bq \bt^{-1} -2 \bq \bt^{-3} +\bq \bt^{-5} -\bq^{-1} \bt^{-3} +\bq^{-1} \bt^{-5} \\
[\ba^{-1}]\colon\ & -\bq^{2} \bt^{-1} +\bq^{2} \bt^{-3} -\bt^{-1} +\bt^{-3} \\
[\ba^{0}]\colon\ & \bq \bt^{-1}
\end{align*}

quotient of the sum, $X/((q - t)(\lambda^2 - t))$ (189 terms):
\begin{align*}
[\ba^{-12}]\colon\ & -\bq^{7} \bt^{-3} -\bq^{7} \bt^{-5} -\bq^{7} \bt^{-7} -\bq^{7} \bt^{-9} -\bq^{5} \bt^{-5} -\bq^{5} \bt^{-7} \\
& \quad -\bq^{5} \bt^{-9} -\bq^{3} \bt^{-7} -\bq^{3} \bt^{-9} -\bq \bt^{-9} +\bq^{-3} \bt^{-11} -\bq^{-5} \bt^{-11} \\
& \quad -\bq^{-5} \bt^{-13} -\bq^{-7} \bt^{-13} -\bq^{-7} \bt^{-15} -\bq^{-9} \bt^{-15} -\bq^{-9} \bt^{-17} -\bq^{-11} \bt^{-17} \\
& \quad -\bq^{-11} \bt^{-19} -\bq^{-13} \bt^{-19} -\bq^{-13} \bt^{-21} -\bq^{-15} \bt^{-21} -\bq^{-15} \bt^{-23} \\
[\ba^{-11}]\colon\ & \bq^{8} \bt^{-9} -\bq^{6} \bt^{-3} -\bq^{6} \bt^{-5} -\bq^{6} \bt^{-7} +\bq^{6} \bt^{-9} -2 \bq^{4} \bt^{-5} \\
& \quad -2 \bq^{4} \bt^{-7} -3 \bq^{2} \bt^{-7} -\bq^{2} \bt^{-9} -2 \bt^{-9} +2 \bq^{-2} \bt^{-9} -\bq^{-2} \bt^{-11} \\
& \quad +4 \bq^{-4} \bt^{-11} +4 \bq^{-6} \bt^{-13} +4 \bq^{-8} \bt^{-15} +3 \bq^{-10} \bt^{-17} +2 \bq^{-12} \bt^{-19} +\bq^{-14} \bt^{-21} \\
[\ba^{-10}]\colon\ & \bq^{9} \bt +\bq^{9} \bt^{-1} +\bq^{9} \bt^{-3} +\bq^{9} \bt^{-5} +\bq^{9} \bt^{-7} +2 \bq^{7} \bt^{-1} \\
& \quad +2 \bq^{7} \bt^{-3} +2 \bq^{7} \bt^{-5} +3 \bq^{7} \bt^{-7} +2 \bq^{5} \bt^{-3} +2 \bq^{5} \bt^{-5} +4 \bq^{5} \bt^{-7} \\
& \quad +\bq^{5} \bt^{-9} +2 \bq^{3} \bt^{-5} +5 \bq^{3} \bt^{-7} +2 \bq^{3} \bt^{-9} +6 \bq \bt^{-7} +4 \bq \bt^{-9} \\
& \quad +4 \bq^{-1} \bt^{-7} +5 \bq^{-1} \bt^{-9} +\bq^{-1} \bt^{-11} +2 \bq^{-3} \bt^{-7} +7 \bq^{-3} \bt^{-9} +3 \bq^{-3} \bt^{-11} \\
& \quad +4 \bq^{-5} \bt^{-9} +7 \bq^{-5} \bt^{-11} +\bq^{-5} \bt^{-13} +4 \bq^{-7} \bt^{-11} +6 \bq^{-7} \bt^{-13} +3 \bq^{-9} \bt^{-13} \\
& \quad +4 \bq^{-9} \bt^{-15} +3 \bq^{-11} \bt^{-15} +3 \bq^{-11} \bt^{-17} +2 \bq^{-13} \bt^{-17} +2 \bq^{-13} \bt^{-19} +\bq^{-15} \bt^{-19} \\
& \quad +\bq^{-15} \bt^{-21} \\
[\ba^{-9}]\colon\ & \bq^{8} \bt +\bq^{8} \bt^{-1} +\bq^{8} \bt^{-3} +\bq^{8} \bt^{-5} -\bq^{8} \bt^{-7} +3 \bq^{6} \bt^{-1} \\
& \quad +3 \bq^{6} \bt^{-3} +3 \bq^{6} \bt^{-5} -\bq^{6} \bt^{-7} +5 \bq^{4} \bt^{-3} +5 \bq^{4} \bt^{-5} +8 \bq^{2} \bt^{-5} \\
& \quad +2 \bq^{2} \bt^{-7} -\bq^{2} \bt^{-9} +6 \bt^{-7} -\bt^{-9} -4 \bq^{-2} \bt^{-7} +\bq^{-2} \bt^{-9} \\
& \quad -\bq^{-2} \bt^{-11} -10 \bq^{-4} \bt^{-9} -\bq^{-4} \bt^{-11} -11 \bq^{-6} \bt^{-11} -\bq^{-6} \bt^{-13} -9 \bq^{-8} \bt^{-13} \\
& \quad -6 \bq^{-10} \bt^{-15} -3 \bq^{-12} \bt^{-17} -\bq^{-14} \bt^{-19} \\
[\ba^{-8}]\colon\ & -\bq^{9} \bt^{3} -\bq^{9} \bt -\bq^{9} \bt^{-1} -\bq^{9} \bt^{-3} -\bq^{9} \bt^{-5} -\bq^{7} \bt \\
& \quad -\bq^{7} \bt^{-1} -\bq^{7} \bt^{-3} -2 \bq^{7} \bt^{-5} -\bq^{5} \bt^{-1} -\bq^{5} \bt^{-3} -3 \bq^{5} \bt^{-5} \\
& \quad -\bq^{5} \bt^{-7} -\bq^{3} \bt^{-3} -5 \bq^{3} \bt^{-5} -3 \bq^{3} \bt^{-7} -6 \bq \bt^{-5} -7 \bq \bt^{-7} \\
& \quad -4 \bq^{-1} \bt^{-5} -11 \bq^{-1} \bt^{-7} -2 \bq^{-1} \bt^{-9} -13 \bq^{-3} \bt^{-7} -6 \bq^{-3} \bt^{-9} -3 \bq^{-5} \bt^{-7} \\
& \quad -12 \bq^{-5} \bt^{-9} -2 \bq^{-5} \bt^{-11} -4 \bq^{-7} \bt^{-9} -8 \bq^{-7} \bt^{-11} -3 \bq^{-9} \bt^{-11} -4 \bq^{-9} \bt^{-13} \\
& \quad -2 \bq^{-11} \bt^{-13} -2 \bq^{-11} \bt^{-15} -\bq^{-13} \bt^{-15} -\bq^{-13} \bt^{-17} \\
[\ba^{-7}]\colon\ & -\bq^{8} \bt^{3} -\bq^{8} \bt -\bq^{8} \bt^{-1} -\bq^{8} \bt^{-3} -2 \bq^{6} \bt -2 \bq^{6} \bt^{-1} \\
& \quad -2 \bq^{6} \bt^{-3} -4 \bq^{4} \bt^{-1} -4 \bq^{4} \bt^{-3} -8 \bq^{2} \bt^{-3} -2 \bq^{2} \bt^{-5} +\bq^{2} \bt^{-7} \\
& \quad -7 \bt^{-5} +2 \bt^{-7} +2 \bq^{-2} \bt^{-5} +\bq^{-2} \bt^{-7} +\bq^{-2} \bt^{-9} +7 \bq^{-4} \bt^{-7} \\
& \quad +2 \bq^{-4} \bt^{-9} +9 \bq^{-6} \bt^{-9} +\bq^{-6} \bt^{-11} +6 \bq^{-8} \bt^{-11} +3 \bq^{-10} \bt^{-13} +\bq^{-12} \bt^{-15} \\
[\ba^{-6}]\colon\ & \bq^{3} \bt^{-3} +\bq^{3} \bt^{-5} +\bq \bt^{-3} +3 \bq \bt^{-5} +6 \bq^{-1} \bt^{-5} +\bq^{-1} \bt^{-7} \\
& \quad +3 \bq^{-3} \bt^{-5} +4 \bq^{-3} \bt^{-7} +6 \bq^{-5} \bt^{-7} +\bq^{-5} \bt^{-9} +\bq^{-7} \bt^{-7} +3 \bq^{-7} \bt^{-9} \\
& \quad +\bq^{-9} \bt^{-9} +\bq^{-9} \bt^{-11} \\
[\ba^{-5}]\colon\ & \bq^{4} \bt +\bq^{4} \bt^{-1} +3 \bq^{2} \bt^{-1} +\bq^{2} \bt^{-3} +3 \bt^{-3} -\bt^{-5} \\
& \quad -\bq^{-4} \bt^{-5} -\bq^{-4} \bt^{-7} -2 \bq^{-6} \bt^{-7} -\bq^{-8} \bt^{-9} \\
[\ba^{-4}]\colon\ & -\bq^{-3} \bt^{-5} \\
[\ba^{-3}]\colon\ & -\bq^{-2} \bt^{-3}
\end{align*}

quotient of $\hat g_{\yone}$:
\begin{align*}
X_{\yone}/((q - t)(\lambda^2 - t)) &= -\ba^{-4} \bq^{-4} \bt^{-4} +\ba^{-5} \bq^{-3} \bt^{-4} +\ba^{-6} \bq^{-4} \bt^{-6}
\end{align*}


\section{Brane--brane and brane--image sectors of the cover}\label{app:cover}
Evidence for \S\ref{sec:split} and \S\ref{sec:kn}. For the two-box trefoil and $T(2,5)$, the sectors of the generating function of the cover split into brane ($Y$) and image ($Y'$), expanded in $G$ with the $t$-side $(t^{-1/2} - t^{1/2})\hat f$ of \cite{KN17} (per $\lambda^\beta$, $\qr = (q/t)^{1/2}$). Brane--brane is $\hat f_{\ytwo}$, $\hat f_{\yoneone}$ (hatted with $M \otimes M$; up to the factor of the uniform normalization it agrees with the mirror of $\hat f$ of \cite{KN17} with $R \leftrightarrow R^T$, \S\ref{sec:split}); brane--image is $\Bim = W_{(1;\bar1)} + \kappa - W_{\yone}^2$ divided by $z_q$ (written $x/z_q$ in the generated tables). The sign classification per block is given in Table \ref{tab:coversplit}: brane--brane is aligned in every block, while $\Bim/z_q$ of brane--image mixes ``aligned'', ``$(\bt^2 - 1)\times$ aligned'' and ``mixed''. For $\Nfix$ with all boundary factors removed (\S\ref{sec:split}; written $N$ in the tables) the only mixed block is the switching block $k = 2m$.
\begin{table}[htbp]
\caption{Signs of the $G$-expansion per $\lambda^\beta$ block for the sectors of the cover.}\label{tab:coversplit}
\begin{center}
\footnotesize
\begin{tabular}{llp{7cm}}
\toprule
knot & sector & sign per $\lambda^\beta$ \\
\midrule
$T(2,3)$ & brane--image $N$ with boundary factors removed & $\lambda^{-12}$: aligned, \allowbreak $\lambda^{-10}$: aligned, \allowbreak $\lambda^{-8}$: aligned, \allowbreak $\lambda^{-6}$: aligned \\
$T(2,3)$ & brane--brane $(t^{-1/2} - t^{1/2})\hat f_{\ytwo}$ & $\lambda^{-10}$: aligned, $\lambda^{-8}$: aligned, $\lambda^{-6}$: aligned, $\lambda^{-4}$: aligned, $\lambda^{-2}$: aligned \\
$T(2,3)$ & brane--brane $(t^{-1/2} - t^{1/2})\hat f_{\yoneone}$ & $\lambda^{-10}$: aligned, $\lambda^{-8}$: aligned, $\lambda^{-6}$: aligned, $\lambda^{-4}$: aligned, $\lambda^{-2}$: aligned \\
$T(2,3)$ & brane--image $(t^{-1/2} - t^{1/2})\,x/z_q$ & $\lambda^{-12}$: aligned, $\lambda^{-10}$: $(\bt^2 - 1)\times$ aligned, $\lambda^{-8}$: \textbf{mixed}, $\lambda^{-6}$: aligned, $\lambda^{-4}$: aligned, $\lambda^{-2}$: $(\bt^2 - 1)\times$ aligned, $\lambda^{0}$: aligned \\
$T(2,5)$ & brane--image $N$ with boundary factors removed & $\lambda^{-20}$: aligned, \allowbreak $\lambda^{-18}$: aligned, \allowbreak $\lambda^{-16}$: aligned, \allowbreak $\lambda^{-14}$: \textbf{mixed}, \allowbreak $\lambda^{-12}$: aligned, \allowbreak $\lambda^{-10}$: aligned, \allowbreak $\lambda^{-8}$: aligned, \allowbreak $\lambda^{-6}$: aligned \\
$T(2,5)$ & brane--brane $(t^{-1/2} - t^{1/2})\hat f_{\ytwo}$ & $\lambda^{-14}$: aligned, $\lambda^{-12}$: aligned, $\lambda^{-10}$: aligned, $\lambda^{-8}$: aligned, $\lambda^{-6}$: aligned \\
$T(2,5)$ & brane--brane $(t^{-1/2} - t^{1/2})\hat f_{\yoneone}$ & $\lambda^{-14}$: aligned, $\lambda^{-12}$: aligned, $\lambda^{-10}$: aligned, $\lambda^{-8}$: aligned, $\lambda^{-6}$: aligned \\
$T(2,5)$ & brane--image $(t^{-1/2} - t^{1/2})\,x/z_q$ & $\lambda^{-20}$: aligned, $\lambda^{-18}$: $(\bt^2 - 1)\times$ aligned, $\lambda^{-16}$: \textbf{mixed}, $\lambda^{-14}$: \textbf{mixed}, $\lambda^{-12}$: \textbf{mixed}, $\lambda^{-10}$: \textbf{mixed}, $\lambda^{-8}$: aligned, $\lambda^{-6}$: $(\bt^2 - 1)\times$ aligned, $\lambda^{-4}$: $(\bt^2 - 1)\times$ aligned, $\lambda^{-2}$: $(\bt^2 - 1)\times$ aligned, $\lambda^{0}$: aligned \\
\bottomrule
\end{tabular}
\end{center}

\end{table}
$T(2,3)$: $\Nfix$ with all boundary factors removed, $\Nfix = (1-t)\Bim/(s^2(t-\lambda^2)(\lambda^2-1)(q-1)(\lambda^2q-1))$ ($\times (qt)^{4/4}$): $G$-expansion (per $\lambda^\beta$, 8 terms)
\begin{align*}
[\lambda^{-12}]\colon\ & -\qr^{-8} \\
[\lambda^{-10}]\colon\ & -G \qr^{-5} -2 \qr^{-4} -\qr^{-6} \\
[\lambda^{-8}]\colon\ & G \qr^{-3} +\qr^{-2} +2 \qr^{-4} \\
[\lambda^{-6}]\colon\ & \qr^{-2}
\end{align*}

$T(2,3)$, brane--brane $(t^{-1/2} - t^{1/2})\hat f_{\ytwo}$: $G$-expansion (per $\lambda^\beta$, 46 terms)
\begin{align*}
[\lambda^{-10}]\colon\ & -G^{2} \qr^{-7} -2 G \qr^{-6} -2 G \qr^{-8} -2 \qr^{-5} -\qr^{-7} -\qr^{-9} \\
[\lambda^{-8}]\colon\ & G^{3} \qr^{-6} +4 G^{2} \qr^{-5} +4 G^{2} \qr^{-7} +6 G \qr^{-4} +9 G \qr^{-6} +5 G \qr^{-8} \\
& \quad +3 \qr^{-3} +7 \qr^{-5} +4 \qr^{-7} +2 \qr^{-9} \\
[\lambda^{-6}]\colon\ & -G^{3} \qr^{-4} -G^{3} \qr^{-6} -3 G^{2} \qr^{-3} -8 G^{2} \qr^{-5} -3 G^{2} \qr^{-7} -3 G \qr^{-2} \\
& \quad -14 G \qr^{-4} -12 G \qr^{-6} -3 G \qr^{-8} -\qr^{-1} -8 \qr^{-3} -9 \qr^{-5} \\
& \quad -5 \qr^{-7} -\qr^{-9} \\
[\lambda^{-4}]\colon\ & G^{3} \qr^{-4} +4 G^{2} \qr^{-3} +4 G^{2} \qr^{-5} +5 G \qr^{-2} +10 G \qr^{-4} +5 G \qr^{-6} \\
& \quad +2 \qr^{-1} +7 \qr^{-3} +5 \qr^{-5} +2 \qr^{-7} \\
[\lambda^{-2}]\colon\ & -G^{2} \qr^{-3} -2 G \qr^{-2} -2 G \qr^{-4} -\qr^{-1} -2 \qr^{-3} -\qr^{-5}
\end{align*}

$T(2,3)$, brane--brane $(t^{-1/2} - t^{1/2})\hat f_{\yoneone}$: $G$-expansion (per $\lambda^\beta$, 27 terms)
\begin{align*}
[\lambda^{-10}]\colon\ & -G \qr^{-6} -\qr^{-5} -\qr^{-7} \\
[\lambda^{-8}]\colon\ & G^{2} \qr^{-5} +3 G \qr^{-4} +3 G \qr^{-6} +2 \qr^{-3} +4 \qr^{-5} +2 \qr^{-7} \\
[\lambda^{-6}]\colon\ & -G^{2} \qr^{-3} -G^{2} \qr^{-5} -2 G \qr^{-2} -6 G \qr^{-4} -2 G \qr^{-6} -\qr^{-1} \\
& \quad -5 \qr^{-3} -5 \qr^{-5} -\qr^{-7} \\
[\lambda^{-4}]\colon\ & G^{2} \qr^{-3} +3 G \qr^{-2} +3 G \qr^{-4} +2 \qr^{-1} +4 \qr^{-3} +2 \qr^{-5} \\
[\lambda^{-2}]\colon\ & -G \qr^{-2} -\qr^{-1} -\qr^{-3}
\end{align*}

$T(2,3)$, brane--image $(t^{-1/2} - t^{1/2})\,\Bim/z_q$: $G$-expansion (per $\lambda^\beta$, 29 terms)
\begin{align*}
[\lambda^{-12}]\colon\ & -\qr^{-9} \\
[\lambda^{-10}]\colon\ & -G \qr^{-6} +G \qr^{-8} -2 \qr^{-5} +2 \qr^{-9} \\
[\lambda^{-8}]\colon\ & G^{2} \qr^{-5} +4 G \qr^{-4} +3 G \qr^{-6} -G \qr^{-8} +3 \qr^{-3} +7 \qr^{-5} \\
& \quad -\qr^{-9} \\
[\lambda^{-6}]\colon\ & -G^{2} \qr^{-3} -G^{2} \qr^{-5} -2 G \qr^{-2} -8 G \qr^{-4} -2 G \qr^{-6} -\qr^{-1} \\
& \quad -7 \qr^{-3} -8 \qr^{-5} \\
[\lambda^{-4}]\colon\ & G^{2} \qr^{-3} +2 G \qr^{-2} +4 G \qr^{-4} +\qr^{-1} +5 \qr^{-3} +3 \qr^{-5} \\
[\lambda^{-2}]\colon\ & \qr^{-1} -\qr^{-3} \\
[\lambda^{0}]\colon\ & -\qr^{-1}
\end{align*}

$T(2,3)$: the same, factorized per $\lambda^\beta$ ($\qr^2 - 1 = \bt^2 - 1$)
\begin{align*}
[\lambda^{-12}]\colon\ & \qr^{-9}\,\bigl(-1\bigr) \\
[\lambda^{-10}]\colon\ & \qr^{-9}\,\bigl(- (\qr^2 - 1)(G \qr + 2 \qr^{2} + 2)\bigr) \\
[\lambda^{-8}]\colon\ & \text{irreducible (as above)} \\
[\lambda^{-6}]\colon\ & \text{irreducible (as above)} \\
[\lambda^{-4}]\colon\ & \text{irreducible (as above)} \\
[\lambda^{-2}]\colon\ & \qr^{-3}\,(\qr^2 - 1) \\
[\lambda^{0}]\colon\ & \qr^{-1}\,\bigl(-1\bigr)
\end{align*}

$T(2,5)$: $\Nfix$ with all boundary factors removed, $\Nfix = (1-t)\Bim/(s^2(t-\lambda^2)(\lambda^2-1)(q-1)(\lambda^2q-1))$ ($\times (qt)^{4/4}$): $G$-expansion (per $\lambda^\beta$, 40 terms)
\begin{align*}
[\lambda^{-20}]\colon\ & -\qr^{-16} \\
[\lambda^{-18}]\colon\ & -G \qr^{-13} -\qr^{-12} -2 \qr^{-14} \\
[\lambda^{-16}]\colon\ & -G^{2} \qr^{-10} -2 G \qr^{-9} -3 G \qr^{-11} -\qr^{-8} -3 \qr^{-10} -2 \qr^{-12} \\
[\lambda^{-14}]\colon\ & -G^{3} \qr^{-7} -4 G^{2} \qr^{-6} -3 G^{2} \qr^{-8} -5 G \qr^{-5} -8 G \qr^{-7} -2 G \qr^{-9} \\
& \quad -2 \qr^{-4} -5 \qr^{-6} -4 \qr^{-8} +\qr^{-10} \\
[\lambda^{-12}]\colon\ & G^{3} \qr^{-5} +3 G^{2} \qr^{-4} +4 G^{2} \qr^{-6} +3 G \qr^{-3} +6 G \qr^{-5} +6 G \qr^{-7} \\
& \quad +\qr^{-2} +2 \qr^{-4} +3 \qr^{-6} +4 \qr^{-8} \\
[\lambda^{-10}]\colon\ & G^{2} \qr^{-4} +2 G \qr^{-3} +3 G \qr^{-5} +\qr^{-2} +2 \qr^{-4} +3 \qr^{-6} \\
[\lambda^{-8}]\colon\ & G \qr^{-3} +\qr^{-2} +2 \qr^{-4} \\
[\lambda^{-6}]\colon\ & \qr^{-2}
\end{align*}

$T(2,5)$, brane--brane $(t^{-1/2} - t^{1/2})\hat f_{\ytwo}$: $G$-expansion (per $\lambda^\beta$, 172 terms)
\begin{align*}
[\lambda^{-14}]\colon\ & -G^{6} \qr^{-11} -6 G^{5} \qr^{-10} -6 G^{5} \qr^{-12} -16 G^{4} \qr^{-9} -26 G^{4} \qr^{-11} -15 G^{4} \qr^{-13} \\
& \quad -24 G^{3} \qr^{-8} -49 G^{3} \qr^{-10} -44 G^{3} \qr^{-12} -20 G^{3} \qr^{-14} -22 G^{2} \qr^{-7} -49 G^{2} \qr^{-9} \\
& \quad -54 G^{2} \qr^{-11} -36 G^{2} \qr^{-13} -15 G^{2} \qr^{-15} -12 G \qr^{-6} -28 G \qr^{-8} -30 G \qr^{-10} \\
& \quad -25 G \qr^{-12} -14 G \qr^{-14} -6 G \qr^{-16} -3 \qr^{-5} -7 \qr^{-7} -8 \qr^{-9} \\
& \quad -5 \qr^{-11} -4 \qr^{-13} -2 \qr^{-15} -\qr^{-17} \\
[\lambda^{-12}]\colon\ & G^{7} \qr^{-10} +8 G^{6} \qr^{-9} +8 G^{6} \qr^{-11} +28 G^{5} \qr^{-8} +50 G^{5} \qr^{-10} +27 G^{5} \qr^{-12} \\
& \quad +56 G^{4} \qr^{-7} +133 G^{4} \qr^{-9} +126 G^{4} \qr^{-11} +50 G^{4} \qr^{-13} +70 G^{3} \qr^{-6} +195 G^{3} \qr^{-8} \\
& \quad +240 G^{3} \qr^{-10} +164 G^{3} \qr^{-12} +55 G^{3} \qr^{-14} +55 G^{2} \qr^{-5} +170 G^{2} \qr^{-7} +237 G^{2} \qr^{-9} \\
& \quad +205 G^{2} \qr^{-11} +116 G^{2} \qr^{-13} +36 G^{2} \qr^{-15} +25 G \qr^{-4} +84 G \qr^{-6} +127 G \qr^{-8} \\
& \quad +117 G \qr^{-10} +82 G \qr^{-12} +42 G \qr^{-14} +13 G \qr^{-16} +5 \qr^{-3} +18 \qr^{-5} \\
& \quad +30 \qr^{-7} +28 \qr^{-9} +19 \qr^{-11} +12 \qr^{-13} +6 \qr^{-15} +2 \qr^{-17} \\
[\lambda^{-10}]\colon\ & -G^{7} \qr^{-8} -G^{7} \qr^{-10} -7 G^{6} \qr^{-7} -16 G^{6} \qr^{-9} -7 G^{6} \qr^{-11} -22 G^{5} \qr^{-6} \\
& \quad -72 G^{5} \qr^{-8} -71 G^{5} \qr^{-10} -21 G^{5} \qr^{-12} -40 G^{4} \qr^{-5} -163 G^{4} \qr^{-7} -228 G^{4} \qr^{-9} \\
& \quad -150 G^{4} \qr^{-11} -35 G^{4} \qr^{-13} -45 G^{3} \qr^{-4} -216 G^{3} \qr^{-6} -368 G^{3} \qr^{-8} -335 G^{3} \qr^{-10} \\
& \quad -175 G^{3} \qr^{-12} -35 G^{3} \qr^{-14} -31 G^{2} \qr^{-3} -172 G^{2} \qr^{-5} -334 G^{2} \qr^{-7} -360 G^{2} \qr^{-9} \\
& \quad -252 G^{2} \qr^{-11} -116 G^{2} \qr^{-13} -21 G^{2} \qr^{-15} -12 G \qr^{-2} -77 G \qr^{-4} -166 G \qr^{-6} \\
& \quad -196 G \qr^{-8} -158 G \qr^{-10} -93 G \qr^{-12} -41 G \qr^{-14} -7 G \qr^{-16} -2 \qr^{-1} \\
& \quad -15 \qr^{-3} -35 \qr^{-5} -46 \qr^{-7} -37 \qr^{-9} -25 \qr^{-11} -13 \qr^{-13} \\
& \quad -6 \qr^{-15} -\qr^{-17} \\
[\lambda^{-8}]\colon\ & G^{7} \qr^{-8} +8 G^{6} \qr^{-7} +8 G^{6} \qr^{-9} +28 G^{5} \qr^{-6} +50 G^{5} \qr^{-8} +27 G^{5} \qr^{-10} \\
& \quad +56 G^{4} \qr^{-5} +133 G^{4} \qr^{-7} +126 G^{4} \qr^{-9} +50 G^{4} \qr^{-11} +69 G^{3} \qr^{-4} +195 G^{3} \qr^{-6} \\
& \quad +241 G^{3} \qr^{-8} +164 G^{3} \qr^{-10} +55 G^{3} \qr^{-12} +52 G^{2} \qr^{-3} +166 G^{2} \qr^{-5} +241 G^{2} \qr^{-7} \\
& \quad +208 G^{2} \qr^{-9} +116 G^{2} \qr^{-11} +36 G^{2} \qr^{-13} +22 G \qr^{-2} +77 G \qr^{-4} +127 G \qr^{-6} \\
& \quad +124 G \qr^{-8} +85 G \qr^{-10} +42 G \qr^{-12} +13 G \qr^{-14} +4 \qr^{-1} +15 \qr^{-3} \\
& \quad +28 \qr^{-5} +30 \qr^{-7} +22 \qr^{-9} +13 \qr^{-11} +6 \qr^{-13} +2 \qr^{-15} \\
[\lambda^{-6}]\colon\ & -G^{6} \qr^{-7} -6 G^{5} \qr^{-6} -6 G^{5} \qr^{-8} -16 G^{4} \qr^{-5} -26 G^{4} \qr^{-7} -15 G^{4} \qr^{-9} \\
& \quad -24 G^{3} \qr^{-4} -49 G^{3} \qr^{-6} -44 G^{3} \qr^{-8} -20 G^{3} \qr^{-10} -21 G^{2} \qr^{-3} -49 G^{2} \qr^{-5} \\
& \quad -55 G^{2} \qr^{-7} -36 G^{2} \qr^{-9} -15 G^{2} \qr^{-11} -10 G \qr^{-2} -25 G \qr^{-4} -33 G \qr^{-6} \\
& \quad -27 G \qr^{-8} -14 G \qr^{-10} -6 G \qr^{-12} -2 \qr^{-1} -5 \qr^{-3} -8 \qr^{-5} \\
& \quad -7 \qr^{-7} -5 \qr^{-9} -2 \qr^{-11} -\qr^{-13}
\end{align*}

$T(2,5)$, brane--brane $(t^{-1/2} - t^{1/2})\hat f_{\yoneone}$: $G$-expansion (per $\lambda^\beta$, 133 terms)
\begin{align*}
[\lambda^{-14}]\colon\ & -G^{5} \qr^{-10} -5 G^{4} \qr^{-9} -5 G^{4} \qr^{-11} -11 G^{3} \qr^{-8} -17 G^{3} \qr^{-10} -10 G^{3} \qr^{-12} \\
& \quad -13 G^{2} \qr^{-7} -25 G^{2} \qr^{-9} -21 G^{2} \qr^{-11} -10 G^{2} \qr^{-13} -8 G \qr^{-6} -18 G \qr^{-8} \\
& \quad -18 G \qr^{-10} -11 G \qr^{-12} -5 G \qr^{-14} -2 \qr^{-5} -6 \qr^{-7} -5 \qr^{-9} \\
& \quad -4 \qr^{-11} -2 \qr^{-13} -\qr^{-15} \\
[\lambda^{-12}]\colon\ & G^{6} \qr^{-9} +7 G^{5} \qr^{-8} +7 G^{5} \qr^{-10} +21 G^{4} \qr^{-7} +37 G^{4} \qr^{-9} +20 G^{4} \qr^{-11} \\
& \quad +35 G^{3} \qr^{-6} +81 G^{3} \qr^{-8} +75 G^{3} \qr^{-10} +30 G^{3} \qr^{-12} +34 G^{2} \qr^{-5} +94 G^{2} \qr^{-7} \\
& \quad +110 G^{2} \qr^{-9} +73 G^{2} \qr^{-11} +25 G^{2} \qr^{-13} +18 G \qr^{-4} +57 G \qr^{-6} +78 G \qr^{-8} \\
& \quad +62 G \qr^{-10} +34 G \qr^{-12} +11 G \qr^{-14} +4 \qr^{-3} +14 \qr^{-5} +23 \qr^{-7} \\
& \quad +19 \qr^{-9} +12 \qr^{-11} +6 \qr^{-13} +2 \qr^{-15} \\
[\lambda^{-10}]\colon\ & -G^{6} \qr^{-7} -G^{6} \qr^{-9} -6 G^{5} \qr^{-6} -14 G^{5} \qr^{-8} -6 G^{5} \qr^{-10} -16 G^{4} \qr^{-5} \\
& \quad -53 G^{4} \qr^{-7} -52 G^{4} \qr^{-9} -15 G^{4} \qr^{-11} -24 G^{3} \qr^{-4} -99 G^{3} \qr^{-6} -135 G^{3} \qr^{-8} \\
& \quad -88 G^{3} \qr^{-10} -20 G^{3} \qr^{-12} -21 G^{2} \qr^{-3} -102 G^{2} \qr^{-5} -171 G^{2} \qr^{-7} -148 G^{2} \qr^{-9} \\
& \quad -77 G^{2} \qr^{-11} -15 G^{2} \qr^{-13} -10 G \qr^{-2} -55 G \qr^{-4} -109 G \qr^{-6} -113 G \qr^{-8} \\
& \quad -73 G \qr^{-10} -34 G \qr^{-12} -6 G \qr^{-14} -2 \qr^{-1} -12 \qr^{-3} -28 \qr^{-5} \\
& \quad -33 \qr^{-7} -25 \qr^{-9} -13 \qr^{-11} -6 \qr^{-13} -\qr^{-15} \\
[\lambda^{-8}]\colon\ & G^{6} \qr^{-7} +7 G^{5} \qr^{-6} +7 G^{5} \qr^{-8} +21 G^{4} \qr^{-5} +37 G^{4} \qr^{-7} +20 G^{4} \qr^{-9} \\
& \quad +35 G^{3} \qr^{-4} +81 G^{3} \qr^{-6} +75 G^{3} \qr^{-8} +30 G^{3} \qr^{-10} +34 G^{2} \qr^{-3} +93 G^{2} \qr^{-5} \\
& \quad +111 G^{2} \qr^{-7} +73 G^{2} \qr^{-9} +25 G^{2} \qr^{-11} +18 G \qr^{-2} +54 G \qr^{-4} +79 G \qr^{-6} \\
& \quad +64 G \qr^{-8} +34 G \qr^{-10} +11 G \qr^{-12} +4 \qr^{-1} +12 \qr^{-3} +22 \qr^{-5} \\
& \quad +21 \qr^{-7} +13 \qr^{-9} +6 \qr^{-11} +2 \qr^{-13} \\
[\lambda^{-6}]\colon\ & -G^{5} \qr^{-6} -5 G^{4} \qr^{-5} -5 G^{4} \qr^{-7} -11 G^{3} \qr^{-4} -17 G^{3} \qr^{-6} -10 G^{3} \qr^{-8} \\
& \quad -13 G^{2} \qr^{-3} -25 G^{2} \qr^{-5} -21 G^{2} \qr^{-7} -10 G^{2} \qr^{-9} -8 G \qr^{-2} -17 G \qr^{-4} \\
& \quad -19 G \qr^{-6} -11 G \qr^{-8} -5 G \qr^{-10} -2 \qr^{-1} -4 \qr^{-3} -6 \qr^{-5} \\
& \quad -5 \qr^{-7} -2 \qr^{-9} -\qr^{-11}
\end{align*}

$T(2,5)$, brane--image $(t^{-1/2} - t^{1/2})\,\Bim/z_q$: $G$-expansion (per $\lambda^\beta$, 93 terms)
\begin{align*}
[\lambda^{-20}]\colon\ & -\qr^{-17} \\
[\lambda^{-18}]\colon\ & -G \qr^{-14} +G \qr^{-16} -\qr^{-13} -\qr^{-15} +2 \qr^{-17} \\
[\lambda^{-16}]\colon\ & -G^{2} \qr^{-11} +G^{2} \qr^{-13} -2 G \qr^{-10} -G \qr^{-12} +4 G \qr^{-14} -G \qr^{-16} \\
& \quad -\qr^{-9} -2 \qr^{-11} +2 \qr^{-13} +2 \qr^{-15} -\qr^{-17} \\
[\lambda^{-14}]\colon\ & -G^{3} \qr^{-8} +G^{3} \qr^{-10} -4 G^{2} \qr^{-7} +5 G^{2} \qr^{-11} -G^{2} \qr^{-13} -5 G \qr^{-6} \\
& \quad -5 G \qr^{-8} +8 G \qr^{-10} +5 G \qr^{-12} -3 G \qr^{-14} -2 \qr^{-5} -4 \qr^{-7} \\
& \quad +\qr^{-9} +7 \qr^{-11} -\qr^{-13} -\qr^{-15} \\
[\lambda^{-12}]\colon\ & G^{4} \qr^{-7} +6 G^{3} \qr^{-6} +5 G^{3} \qr^{-8} -G^{3} \qr^{-10} +12 G^{2} \qr^{-5} +23 G^{2} \qr^{-7} \\
& \quad +4 G^{2} \qr^{-9} -4 G^{2} \qr^{-11} +10 G \qr^{-4} +29 G \qr^{-6} +23 G \qr^{-8} -8 G \qr^{-10} \\
& \quad -4 G \qr^{-12} +3 \qr^{-3} +11 \qr^{-5} +15 \qr^{-7} +4 \qr^{-9} -8 \qr^{-11} \\
[\lambda^{-10}]\colon\ & -G^{4} \qr^{-5} -G^{4} \qr^{-7} -4 G^{3} \qr^{-4} -12 G^{3} \qr^{-6} -4 G^{3} \qr^{-8} -6 G^{2} \qr^{-3} \\
& \quad -28 G^{2} \qr^{-5} -32 G^{2} \qr^{-7} -4 G^{2} \qr^{-9} -4 G \qr^{-2} -24 G \qr^{-4} -44 G \qr^{-6} \\
& \quad -30 G \qr^{-8} +2 G \qr^{-10} -\qr^{-1} -7 \qr^{-3} -17 \qr^{-5} -19 \qr^{-7} \\
& \quad -7 \qr^{-9} +3 \qr^{-11} \\
[\lambda^{-8}]\colon\ & G^{4} \qr^{-5} +4 G^{3} \qr^{-4} +6 G^{3} \qr^{-6} +6 G^{2} \qr^{-3} +16 G^{2} \qr^{-5} +13 G^{2} \qr^{-7} \\
& \quad +4 G \qr^{-2} +14 G \qr^{-4} +20 G \qr^{-6} +12 G \qr^{-8} +\qr^{-1} +4 \qr^{-3} \\
& \quad +8 \qr^{-5} +9 \qr^{-7} +3 \qr^{-9} \\
[\lambda^{-6}]\colon\ & \qr^{-5} -\qr^{-7} \\
[\lambda^{-4}]\colon\ & \qr^{-3} -\qr^{-5} \\
[\lambda^{-2}]\colon\ & \qr^{-1} -\qr^{-3} \\
[\lambda^{0}]\colon\ & -\qr^{-1}
\end{align*}

$T(2,5)$: the same, factorized per $\lambda^\beta$ ($\qr^2 - 1 = \bt^2 - 1$)
\begin{align*}
[\lambda^{-20}]\colon\ & \qr^{-17}\,\bigl(-1\bigr) \\
[\lambda^{-18}]\colon\ & \qr^{-17}\,\bigl(- (\qr^2 - 1)(G \qr + \qr^{2} + 2)\bigr) \\
[\lambda^{-16}]\colon\ & \qr^{-17}\,\bigl(- (\qr^2 - 1)(G \qr + \qr^{2} + 1)(G \qr^{3} + \qr^{4} + 2 \qr^{2} - 1)\bigr) \\
[\lambda^{-14}]\colon\ & \qr^{-15}\,\bigl(- (\qr^2 - 1) \\
& \quad (G^{3} \qr^{5} + 4 G^{2} \qr^{6} + 4 G^{2} \qr^{4} - G^{2} \qr^{2} \\
& \quad  + 5 G \qr^{7} + 10 G \qr^{5} + 2 G \qr^{3} - 3 G \qr + 2 \qr^{8} + 6 \qr^{6} \\
& \quad  + 5 \qr^{4} - 2 \qr^{2} - 1)\bigr) \\
[\lambda^{-12}]\colon\ & \text{irreducible (as above)} \\
[\lambda^{-10}]\colon\ & \text{irreducible (as above)} \\
[\lambda^{-8}]\colon\ & \text{irreducible (as above)} \\
[\lambda^{-6}]\colon\ & \qr^{-7}\,(\qr^2 - 1) \\
[\lambda^{-4}]\colon\ & \qr^{-5}\,(\qr^2 - 1) \\
[\lambda^{-2}]\colon\ & \qr^{-3}\,(\qr^2 - 1) \\
[\lambda^{0}]\colon\ & \qr^{-1}\,\bigl(-1\bigr)
\end{align*}


\section{The sector brane 2, image 2}\label{app:n2}
We summarize the results for the fixed sector $n = 2$ (brane 2, image 2) of the trefoil. $E$ is determined up to $\lambda^{-12}$ and undetermined from $\lambda^{-14}$ on. Notation: the rows of the mirror images of the fundamental and two-box representations are written $R_{\yone} = a_2\lambda^{-2} + a_4\lambda^{-4}$, $R_{\ytwo} = b_4\lambda^{-4} + b_6\lambda^{-6} + b_8\lambda^{-8}$, $R_{\yoneone} = c_4\lambda^{-4} + \cdots$, those of the adjoint $g^{\mathrm{adj}}_j$, and those of the composites $\Gamma_{\alpha\beta}$ (Result \ref{cr:composite}); the ``background'' is the part of the fixed sector that survives when all Wilson loop data are set to their unknot values, and its ``cells'' are the monomials generating it as a Fock space.
% [P30, P31, P38, P39]
\begin{result}[rows of $n = 2$ and the $R$ dependence]\label{cr:n2rows}
The fixed sector $\Nfix_R$ of $n = 2$ for the trefoil ($R = (2), (1,1)$) can be written explicitly as a polynomial (of degree at most four) in the rows $a_j$ ($a_2 = (1+qt)/q$, $a_4 = -t/q$), $b_j$, $c_j$, $g^{\mathrm{adj}}_j$ and $\Gamma_{\alpha\beta}$; with the actual rows it agrees with the computed $\Nfix_R$ in all blocks up to $\lambda^{-20}$.
\begin{enumerate}
\item \textbf{$\sigma$-covariance.} Under the reweighting $\sigma$ of the two-box rows and of $\Gamma$, $b_j \mapsto qb_j$, $c_j \mapsto c_j/t$, $\Gamma_{\alpha\beta} \mapsto w_\alpha w_\beta \Gamma_{\alpha\beta}$ ($w_{\ytwo} = q$, $w_{\yoneone} = 1/t$), the difference $\Nfix_{\ytwo}\circ\sigma - \Nfix_{\yoneone}$ is $(q-t)\times$ a polynomial in the adjoint-sector symbols $a_2, a_4, g^{\mathrm{adj}}_j$ and the background alone (up to $\lambda^{-20}$). The $R$ dependence thus splits into a reweighting of the two-box rows and a tower ``two vacuum states $(q-t)\otimes$ adjoint''.
\item \textbf{The square at $\lambda^{-12}$.} The $\sigma$-covariant part is a perfect square $-(t^{3/2}/q^{1/2})\hat Y_R^2$ for both $R$, with $\hat Y_{\yoneone} = (1+q)(1+t)(1+qt)/q^3$ and $\hat Y_{\ytwo} = [(1+q)(1+t)(1+qt)^2 + qt(q-t)]/q^4$ (10 monomials, positive, irreducible, not a product). The correction is $(q-t)[(t/q^4)a_2^2 + \eta]$ with $\eta = [tV'^2a_2^2/q^3 - 2a_2\tilde\Sigma + (V'-t)g^{\mathrm{adj}}_6]/q^3$ ($V' = 1+q+qt$, $\tilde\Sigma = (1+t+qt)/q$; $\eta = 0$ for $R = (1,1)$). Distinguishing the brane and image rows by different symbols, this square is literally the tensor product $X^bX^i$ of the brane $X$ and the image $X$ ($X$ a linear combination of $a_2^2, b_4, c_4$).
\item \textbf{$\hat Y_R$ is a single-brane quantity.} The two-box cumulants of a single brane ($\log$ of $W_{\yone}, W_{\ytwo}, W_{\yoneone}$ with the Adams term $f_{\ytwo} = Z_{\ytwo} - \tfrac12\psi^2 Z_{\yone}$) projected onto the Schur $R$, $\Omega^{(0)}_R(\lambda) = \sum_\rho \chi_R(\rho)f_\rho/z_\rho$ (with $z_\rho$ the order of the centralizer of the cycle type $\rho$), equal $\mathrm{fac}\cdot[\Omega'_2\lambda^{-2} + \Omega'_4\lambda^{-4} + \cdots]$ with $\mathrm{fac} = -q^{1/2}t/(t-1)$, and exactly $\Omega'_4 = \hat Y_R$ for both $R$. The $qt(q-t)$ in $\hat Y_{\ytwo}$ comes from the Adams term. For $R = (1,1)$ the correction is also the square of $\Omega'_2 = -a_2/q$: $\Nfix_{12}(1,1) = -(t^{3/2}/q^{1/2})[\Omega_4'^2 + (q-t)(t/q^2)\Omega_2'^2]$.
\end{enumerate}
\end{result}
% [P40]
\begin{result}[generating function of the adjoint tower]\label{cr:tower}
Let $\xi = \lambda^{-2}$, $V = 1+t+qt$, $V' = 1+q+qt$. The kernels in $\Nfix_{\ytwo}\circ\sigma - \Nfix_{\yoneone} = -(t^{3/2}/q^{1/2})(q-t)[R_{\yone}(\xi)^2\kappa_R(\xi) + \Gamma_{\mathrm{adj}}(\xi)\kappa_\Gamma(\xi) + \beta(\xi)]$ are, in closed form,
\begingroup\small
\begin{align*}
\kappa_\Gamma(\xi) &= \frac{q^2\xi^3[\,q^2(V'-t) - t^2\xi(V-q)\,]}{(t\xi-q)(t^3\xi-q)(t^2\xi-q^2)(t\xi-q^3)},\qquad
\beta(\xi) = \frac{-q^2\xi^4}{(\xi-q)(t\xi-q)(t^2\xi-q)(t\xi-q^2)(t\xi-1)},\\
\kappa_R(\xi) &= \frac{-q^2\xi^3\,P_5(\xi)}{(\xi-q)(t\xi-q)(t^2\xi-q)(t^3\xi-q)(t\xi-q^2)(t^2\xi-q^2)(t\xi-q^3)(t\xi-1)}
\end{align*}
\endgroup
with $P_5$ an irreducible polynomial of degree five in $\xi$. They were obtained by running the computation with the two-box rows and $\Gamma$ set to zero as rational functions in $\lambda$ (agreeing with the series in symbolic blocks up to $\lambda^{-20}$). All kernels are ``two vacuum states $(q-t)\times$ Fock kernels of cells'', the cells being the six background cells ($t, 1/q, t/q, t^2/q, t/q^2$; $t^2/q^2$) plus $t^3/q, t/q^3$ specific to $\Gamma_{\mathrm{adj}}$. The series coefficients of $\kappa_\Gamma$ are of staircase type, $t^nT_nU_n/q^{3(n+1)}$ with $T_n = \sum_{0\le j\le i\le n}q^it^j$ and $U_n$ a polynomial of degree $n$ in $q, t$ (not reproduced here).
\end{result}
% [P32, P33, P34, P41]
\begin{result}[$E$ at $\lambda^{-12}$ for $n = 2$: full exchange]\label{cr:n2E}
At $\lambda^{-12}$ the $E$ of $n = 2$ is, for both $R$, ``signs of the background cells plus the exchange $\psi^2$ on all squares and products of rows of the knot part'', and satisfies (E1)(E2):
\begin{gather*}
E_{12}(1,1) = q^{-13/2}t^{3/2}\bigl[\psi^2\hat Y_{\yoneone} + (q-t)(t/q^2)\,\psi^2\Omega'_2\bigr] - B_{12}(1,1),\\
B_{12}(1,1) = -q^{-13/2}t^{3/2}(q-t)qt^2 .
\end{gather*}
The extra $(q-t)qt^2$ of an earlier interpretation (correction $t\,\Sym^2(1+qt)$) was the $E$ of the background ($E_B = -B$). For $R = (2)$, besides the square $\psi^2\hat Y_{\ytwo}$, each piece of the correction is fixed including its sign ($(t/q^4)a_2^2$, $tV'^2a_2^2/q^6$ by $\psi^2$; the cross term $-2a_2\tilde\Sigma$ is $0$; the top row of $\Gamma_{\mathrm{adj}}$ by the same exchange as for $n = 1$, $\gamma_1 = (2\rho_0\rho_1 - \rho_1^2)/q \mapsto -\rho_1^2/q$; its kernel $V' - t$ becomes $(1+q)(1-t)$ by the cell sign $(-1)^{\deg_q\deg_t}$), one of 16 sign choices for the kernel. This $E$ lies outside the earlier family with $|e|\le|R|$ imposed monomial by monomial on the correction. The block signs are unified as $(-1)^{k/2+1}$ ($k$ the exponent in $\lambda^{-k}$) across the $n = 1$ switching, $\lambda^{-10}$ of $R = (2)$ and $\lambda^{-12}$ of both $R$, corresponding to flipping the sign of cell 1 in $F^\varepsilon$ of $n = 1$. By the same rule the second term $V - q$ in the numerator of $\kappa_\Gamma$ becomes $(1-q)(1+t)$ (a conjecture).
\end{result}
% [P35, P36, P37, P42]
\begin{result}[obstruction from $\lambda^{-14}$ on]\label{cr:n2negative}
With the $E$ of $\lambda^{-8}$--$\lambda^{-12}$ fixed, the integer program at $\lambda^{-14}$ determines only 2--3 of 113 unknowns even under (E1)+(E2) (the minimal number of exchanged pairs is 48 under (E1), 936 under (E1)+(E2)). A construction taking the trace of the exchange monomial by monomial in the symbolic blocks does not give an $E$ for any assignment of brane/image labels (either $R$-independent or not Laurent): the Laurent property of $\Nfix$ arises from cancellations between monomials of different content, so the trace is defined only on the brane object $\hat Y_R = \Omega'_4$ in the Schur basis. Pure brane$\times$image products cancel in the mixed cumulant (rank 0), and the square $\hat Y_R^2$ is carried by the connected kernels (dimension correction of the composite, $c_1$, $\kappa$). At $\lambda^{-14}$ the $\sigma$-covariant part is not divisible by $\hat Y_R$ and cannot be written as a product series at the level of $\Omega^{(0)}$.
\end{result}

\section{Unrefined check}\label{app:unref}
The $\hat g$ of the trefoil (standard framing) of \cite{Marino09} computed with our code. $\hat g_{\yoneone}$ agrees coefficient by coefficient with the printed (4.8) of \cite{Marino09}; $\hat g_{\ytwo}$ agrees except for the $z^9$ term (Result \ref{cr:marino}).
$\hat g_{\ytwo}$ (by powers of $z = t_M - t_M^{-1}$; $\lambda$ is the $\nu$ of \cite{Marino09}):
\begin{align*}
[z^{0}]\colon\ & -21 \lambda^{11} +79 \lambda^{9} -111 \lambda^{7} +69 \lambda^{5} -16 \lambda^{3} \\
[z]\colon\ & 21 \lambda^{12} -63 \lambda^{10} +63 \lambda^{8} -21 \lambda^{6} \\
[z^{2}]\colon\ & -70 \lambda^{11} +251 \lambda^{9} -307 \lambda^{7} +146 \lambda^{5} -20 \lambda^{3} \\
[z^{3}]\colon\ & 70 \lambda^{12} -231 \lambda^{10} +231 \lambda^{8} -70 \lambda^{6} \\
[z^{4}]\colon\ & -84 \lambda^{11} +330 \lambda^{9} -366 \lambda^{7} +128 \lambda^{5} -8 \lambda^{3} \\
[z^{5}]\colon\ & 84 \lambda^{12} -322 \lambda^{10} +322 \lambda^{8} -84 \lambda^{6} \\
[z^{6}]\colon\ & -45 \lambda^{11} +220 \lambda^{9} -230 \lambda^{7} +56 \lambda^{5} -\lambda^{3} \\
[z^{7}]\colon\ & 45 \lambda^{12} -219 \lambda^{10} +219 \lambda^{8} -45 \lambda^{6} \\
[z^{8}]\colon\ & -11 \lambda^{11} +78 \lambda^{9} -79 \lambda^{7} +12 \lambda^{5} \\
[z^{9}]\colon\ & 11 \lambda^{12} -78 \lambda^{10} +78 \lambda^{8} -11 \lambda^{6} \\
[z^{10}]\colon\ & -\lambda^{11} +14 \lambda^{9} -14 \lambda^{7} +\lambda^{5} \\
[z^{11}]\colon\ & \lambda^{12} -14 \lambda^{10} +14 \lambda^{8} -\lambda^{6} \\
[z^{12}]\colon\ & \lambda^{9} -\lambda^{7} \\
[z^{13}]\colon\ & -\lambda^{10} +\lambda^{8}
\end{align*}

$\hat g_{\yoneone}$ (by powers of $z = t_M - t_M^{-1}$; $\lambda$ is the $\nu$ of \cite{Marino09}):
\begin{align*}
[z^{0}]\colon\ & -15 \lambda^{11} +53 \lambda^{9} -69 \lambda^{7} +39 \lambda^{5} -8 \lambda^{3} \\
[z]\colon\ & 15 \lambda^{12} -45 \lambda^{10} +45 \lambda^{8} -15 \lambda^{6} \\
[z^{2}]\colon\ & -35 \lambda^{11} +126 \lambda^{9} -146 \lambda^{7} +61 \lambda^{5} -6 \lambda^{3} \\
[z^{3}]\colon\ & 35 \lambda^{12} -120 \lambda^{10} +120 \lambda^{8} -35 \lambda^{6} \\
[z^{4}]\colon\ & -28 \lambda^{11} +120 \lambda^{9} -128 \lambda^{7} +37 \lambda^{5} -\lambda^{3} \\
[z^{5}]\colon\ & 28 \lambda^{12} -119 \lambda^{10} +119 \lambda^{8} -28 \lambda^{6} \\
[z^{6}]\colon\ & -9 \lambda^{11} +55 \lambda^{9} -56 \lambda^{7} +10 \lambda^{5} \\
[z^{7}]\colon\ & 9 \lambda^{12} -55 \lambda^{10} +55 \lambda^{8} -9 \lambda^{6} \\
[z^{8}]\colon\ & -\lambda^{11} +12 \lambda^{9} -12 \lambda^{7} +\lambda^{5} \\
[z^{9}]\colon\ & \lambda^{12} -12 \lambda^{10} +12 \lambda^{8} -\lambda^{6} \\
[z^{10}]\colon\ & \lambda^{9} -\lambda^{7} \\
[z^{11}]\colon\ & -\lambda^{10} +\lambda^{8}
\end{align*}

